\section{Introduction}\label{sec:introduction}

The word problem asks whether a word in a fixed finite generating set
and its inverses represents the identity. A finitely generated group has
\emph{decidable word problem} if an algorithm answers this question for
every word. Embedding theorems connect this algorithmic condition with
the algebraic structure of an overgroup. Higman's embedding theorem
characterizes the finitely generated subgroups of finitely presented
groups as the recursively presented groups: those admitting a
recursively enumerable set of defining relations \citep{Higman1961}.
Boone and Higman obtained the finer characterization
\citep[Theorem~I]{BH74}: a finitely generated group has decidable word
problem if and only if it embeds in a simple group which in turn embeds
in a finitely presented group. They asked whether the simple overgroup
itself can always be finitely presented \citep[p.~43]{BH74}.
This is the Boone--Higman conjecture. Thompson subsequently showed that
the intermediate simple group can be chosen finitely generated with
decidable word problem \citep{Tho80}.

We consider the strengthening that asks for finiteness in every
homotopical dimension. A group \(H\) has \emph{type \(F_n\)} if it admits a
classifying space \(K(H,1)\) with finite \(n\)-skeleton, and has
\emph{type \(\Finfty\)} if it admits such a space with finitely many
cells in each dimension. Here a classifying space is a connected CW
complex with fundamental group \(H\) and contractible universal cover.
Type \(F_1\) is finite generation and type \(F_2\) is finite
presentation. The question of a simple overgroup of type \(\Finfty\)
appears in
the discussion following Question~5.6 of \citet{BBMZsurvey};
Question~5.6 itself asks about overgroups of type \(\Finfty\) for
arbitrary finitely presented groups.

\begin{theorem}\label{thm:main}
Every finitely generated group \(G\) with decidable word problem admits
an injective homomorphism into a nontrivial simple group \(H\) of type
\(\Finfty\).
\end{theorem}

The input need not be finitely presented, and its word-problem
algorithm may have arbitrary running time. The classifying space of
\(H\) may have infinitely many dimensions; type \(\Finfty\) imposes
neither torsion-freeness nor finite cohomological dimension.
The theorem answers the higher-finiteness question positively and in
particular implies the Boone--Higman conjecture.

The classical converse gives an algebraic characterization here as
well \citep{BH74}. A nontrivial finitely presented simple group has
decidable word
problem: in parallel, enumerate consequences of its defining relations
and consequences after adding the input word as a relation. A trivial
input word is detected by the first search; a nontrivial one normally
generates the simple group, so the second search eventually proves all
generators trivial. These outcomes are exclusive because the original
group is nontrivial. Restricting this algorithm along an embedding
solves the word problem in any finitely generated subgroup. Thus the
groups in Theorem~\ref{thm:main} are exactly the finitely generated
subgroups of simple groups of type \(\Finfty\).

\subsection{Earlier embeddings and the higher-finiteness obstacle}

Thompson-type groups have long supplied simple groups with strong
finiteness properties. Brown established type \(\Finfty\) for the
Thompson--Higman families \citep[Section~4]{Brown1987}.
More recently, Belk and Zaremsky constructed a single simple group of
type \(\Finfty\) containing all finitely generated right-angled Artin
groups \citep[Corollary~E]{BelkZaremsky2022}; these are the groups
presented by generators indexed by the vertices of a finite graph,
with commutation relations for its edges. Belk, Hyde and Matucci
constructed a simple group of type \(\Finfty\) containing every
countable abelian group \citep[Theorems~4.4--4.6]{BHMgerms}.
Such results give substantial classes of examples without a general
construction from a word-problem algorithm.

Other recent advances concern the finite-presentation conclusion of
the Boone--Higman conjecture. Belk, Bleak, Matucci and Zaremsky proved
it for hyperbolic groups and contracting self-similar groups
\citep[Theorem~A and Corollary~D]{BBMZhyperbolic}.
Belk, Fournier-Facio, Hyde and Zaremsky proved it for
\(\operatorname{Aut}(F_n)\) and mapping class groups of orientable
punctured surfaces of finite type
\citep[Theorem~A and Corollary~B]{BFFHZ}.
Here \(F_n\) denotes the free group of rank \(n\).
These conclusions do not by themselves give simple overgroups of type
\(\Finfty\). Higher finiteness requires control of additional
relations among relations in every dimension.

Our final step uses the twisted Brin--Thompson construction of
\citet{BelkZaremsky2022}. Its input is a faithful action of a group
\(\Lambda\) on a countably infinite set \(S\). The action must be
\emph{oligomorphic}, meaning that it has finitely many orbits on finite
subsets of each fixed cardinality. If \(\Lambda\) and all its finite-set
stabilizers have type \(\Finfty\), their theorem gives a simple group
\(SV_\Lambda\) of type \(\Finfty\) containing \(\Lambda\).
Consequently the problem is to embed \(G\) in a group equipped with
such an action, including the finiteness properties of its stabilizers.
We construct a \emph{highly transitive} action: any two ordered tuples
of distinct points of the same finite length lie in the same orbit.
This supplies the orbit condition. We prove type \(\Finfty\) for the
pointwise stabilizers; the corresponding setwise stabilizers are
finite extensions and inherit the property.

The stabilizer requirement is a genuine restriction on the choice of
acting group. Fournier-Facio, Kropholler, Lyman and Zaremsky prove that
an oligomorphic action on an infinite set of a group of finite virtual
cohomological dimension, or of a countable group linear over a field,
always has a finite-set stabilizer which is not of type \(\Finfty\)
\citep[Theorem~1.1]{FFKLZstabilizers}.
The construction below therefore builds the overgroup and its
stabilizers together; it does not ask the original group to possess
an action with these properties.

\subsection{Constructing the action and proving finiteness}

The action begins with a finitely presented ring \(B\), the elementary
matrix group \(Q=\E_7(B)\), and its action on the column set \(X=B^7\).
Here \(\E_r(B)\) is generated by the off-diagonal elementary matrices.
We construct compatible injections
\[
 \delta:Q\longrightarrow Q,\qquad \sigma:X\longrightarrow X.
\]
At each new stage, matrices can permute any finite collection of the
compressed columns \(\sigma(X)\). The direct limits of these two
systems, with the stage shift adjoined to the group, give the action
\(\Lambda\curvearrowright S\). This makes high transitivity a local
matrix calculation. The same construction embeds \(G\) into \(Q\).

The ring has to encode more than the given group. Its units must also
encode prescribed enlargements of the action of \(\E_7(B)\) on
\(B^7\), so the objects being encoded already involve the ring being
constructed. We resolve this self-reference by a length induction.
A finite compiler from the companion article
\emph{Finite algebraic envelopes and the Boone--Higman conjecture}
\citep{FoundationsCompanion} gives equality calculations at each finite
bound when its stated hypotheses hold. The encoding ensures that every recursive call uses
a strictly smaller bound. Kleene's recursion theorem first supplies the
program, and the length induction then proves its totality and
compatibility with multiplication. The companion also supplies a
binary corner construction and a Steinberg-kernel annihilation theorem;
their exact statements appear where they are used.

Finiteness is supplied by a separate ascending-torus criterion.
For an injective endomorphism \(f:U\to U\), write
\[
 T(U,f)=\langle U,t\mid t^{-1}ut=f(u)\ \text{for }u\in U\rangle.
\]
A factorization of \(f\) through a finitely presented group gives a
finite presentation of \(T(U,f)\). To obtain higher finiteness, we use
two finite systems of homomorphisms whose images commute in specified
pairs. Their product identities come from orthogonal idempotents over
\(\F_2\) and \(\F_3\). After lower-dimensional terms have been killed,
commuting products act additively on homology. The two systems then
force the same class to be annihilated by both \(2\) and \(3\).

The essential point is uniformity. We use simplices whose vertices are
elements of \(U\), and regard two simplices as having the same type
when they differ by a simultaneous left translation. One iterate and
one finite set of output types work for every cycle using a given
finite set of input types. Chain lengths and coefficients may be
arbitrarily large. This control allows fillings with coefficients in
products of copies of the group ring, and the Bieri--Eckmann finiteness
method then gives finite skeleta in every dimension. The result is
Theorem~\ref{thm:torus}.

The acting group and each finite-point stabilizer are themselves
ascending tori to which this same criterion applies. For stabilizers,
the defining endomorphism factors through a finitely presented
Steinberg group. The second factor sends the \emph{entire} Steinberg
group into the stabilizer, not merely the image of the first factor.
The commuting diagrams preserve the fixed points as well. These two
properties establish type \(\Finfty\) for all stabilizers and complete
the input to the twisted Brin--Thompson theorem.

\subsection{Organization and conventions}

Section~\ref{sec:homological} proves the ascending-torus criterion.
Section~\ref{sec:encoding} constructs the ring and the encoded actions.
Section~\ref{sec:stage} constructs \(\Lambda\curvearrowright S\) and
proves the stabilizer theorem. Section~\ref{sec:simple} applies the
twisted Brin--Thompson theorem and proves Theorem~\ref{thm:main}.

All rings used for matrices are unital and associative. Ring embeddings
and endomorphisms need not preserve the identity when this is expressly
stated. All homology is integral. Group actions are on the left;
coefficient modules for group homology are right modules.
